% !TEX program = pdflatex
\documentclass[3p,11pt]{elsarticle}
\usepackage{amsmath,amssymb}
\usepackage{array}
\usepackage{booktabs}
\usepackage[tableposition=top]{caption}
\usepackage{enumitem}
\usepackage{tikz}
\usepackage{pgfplots}
\pgfplotsset{compat=1.18}
\usepackage{url}
\urlstyle{same}
\usepackage[hidelinks]{hyperref}
\emergencystretch=3em
% first-page footer: the page number alone
\makeatletter
\def\ps@pprintTitle{\let\@oddhead\@empty\let\@evenhead\@empty
  \def\@oddfoot{\hfil\thepage\hfil}\let\@evenfoot\@oddfoot}
\makeatother

% ---- notation ---------------------------------------------------------------
\newcommand{\um}{\,$\mu$m}
\newcommand{\kms}{km\,s$^{-1}$}
\newcommand{\flux}{erg\,cm$^{-2}$\,s$^{-1}$}
\newcommand{\mA}{m\AA}
\newcommand{\sci}[2]{#1\times10^{#2}}
\newcommand{\ion}[2]{#1\,\textsc{#2}}
\newcommand{\nuc}[2]{$^{#1}$#2}
\newcommand{\lxl}{$L_X/L_\mathrm{bol}$}

\begin{document}

\begin{frontmatter}

\title{Persistent narrowband radio lines below the ionospheric cutoff: 27 years of Wind and five years of STEREO}

\author{Charles C. Norton}

\begin{abstract}
Below about $10$\,MHz the ionosphere reflects radio waves, so a transmitter at those
frequencies can be heard only from space. We search every one-minute spectrum of the
Wind/WAVES receivers from 1995 to 2021, $20$\,kHz to $13.8$\,MHz, and every spectrum of
STEREO-A in 2008, 2012 and 2020 and of STEREO-B in 2010 and 2012, $125$\,kHz to
$16$\,MHz, for lines that persist: a channel five robust deviations above its
neighbours in more than one per cent of the minutes on thirty days or more. Every
persistent line has a known source. The solar wind's plasma line, Earth's auroral
kilometric radiation and the terrestrial broadcast, maritime, aeronautical and amateur
services account for most, and terrestrial leakage is recognised by its day-to-day
correlation with the broadcast lines of the 41\,m band, which reaches the spacecraft
only through a sunlit ionosphere. The rest are made on the spacecraft: Wind's lines at
$8025$ and $1175$\,kHz, which do not follow the broadcast lines, and lines that one
STEREO spacecraft holds five times as often as Wind or the other spacecraft, among them
$10425$\,kHz in both STEREO receivers, whose occupancy in STEREO-A rose from $0.08$ per
cent in 2008 to $4.7$ per cent in 2020. No persistent line remains outside the known
sources in any of the $32$ spacecraft-years. The channels are $3$, $20$ and $50$\,kHz
wide, so the search is for narrow lines only in the sense that a line fills one
channel; a $5\sigma$ line in Wind's RAD2 carries $2\times10^{-16}$\,W\,m$^{-2}$, the
output of a $6\times10^{7}$\,W isotropic source at $1$\,au or of $2\times10^{18}$\,W at
$1$\,pc, so the search is one of the Solar System and of the strongest conceivable
transmitters at the nearest stars; Wind's RAD1 measures only $32$ to $38$ of its $256$
channels directly, which alone carry its search, and $64$ per cent of RAD2's band and
$52$ per cent of STEREO's lie inside the terrestrial service bands, where a line is
assigned to the service and not counted.
\end{abstract}

\begin{keyword}
Technosignatures \sep Radio lines: general \sep Solar wind \sep Space vehicles: instruments \sep
Interplanetary medium
\end{keyword}

\end{frontmatter}

% =============================================================================
\section{Introduction}\label{sec:intro}
% =============================================================================

Radio searches for technosignatures listen from the ground, where the ionosphere
reflects everything below its plasma frequency of a few to about ten megahertz, and
searches at low frequencies, such as one with the Murchison Widefield Array near
$100$\,MHz~\cite{Tingay2016}, observe well above it. The band from tens
of kilohertz to ten megahertz is open only to receivers in space. The interstellar plasma
frequency is a few kilohertz, but free-free absorption~\cite{RybickiLightman1979} in the
warm, partly ionised clouds around the Sun, with $n_e \approx 0.1$\,cm$^{-3}$ at about
$7000$\,K~\cite{RedfieldLinsky2008}, sets the reach: across $5$ to $10$\,pc of such gas
the optical depth is about $0.1$ at $1$\,MHz, $1$ near $300$\,kHz and $10$ at $100$\,kHz,
so nearby stars are heard through it above a few hundred kilohertz, depending on the
sightline, and below about $100$\,kHz only sources within the heliosphere or across the
shortest paths through warm gas. Spacecraft built to study the solar wind and the
magnetosphere have recorded that band continuously for three decades: the WAVES
receivers of Wind~\cite{Bougeret1995}, since 2004 at the Sun-Earth L1 point, and those of
the twin STEREO spacecraft~\cite{Bougeret2008}, which drifted ahead of and behind Earth
on its orbit from 2006.

Their spectra are crowded with natural lines and with terrestrial ones. Earth's auroral
kilometric radiation reaches far beyond the magnetosphere~\cite{Gurnett1974}, and
terrestrial transmitters leak through the ionosphere where it is sunlit, as the Radio
Astronomy Explorer found~\cite{Herman1973}. A narrowband transmitter elsewhere would
add a line that persists in one channel, whose behaviour from minute to minute and day to
day separates it from these. This paper searches the full Wind and STEREO records for
such lines.

\subsection{Contributions}
\begin{enumerate}[leftmargin=*,itemsep=1pt]
\item A search of $27$ years of Wind/WAVES and $5$ years of STEREO/WAVES spectra for
persistent narrowband lines below the ionospheric cutoff.
\item Two tests that separate terrestrial and spacecraft lines from the rest: a line's
day-to-day correlation with the broadcast lines, and its occupancy in the other
spacecraft, which share the sky but not the receiver.
\end{enumerate}

% =============================================================================
\section{Data and method}\label{sec:method}
% =============================================================================

\emph{Data.} Wind's RAD1 receiver covers $20$ to $1040$\,kHz in $256$ channels $4$\,kHz
apart and $3$\,kHz wide, and RAD2 $1.075$ to $13.825$\,MHz in $256$ channels $50$\,kHz
apart and $20$\,kHz wide~\cite{Bougeret1995}; CDAWeb's WI\_H1\_WAV files give one-minute
averages of both for every day from 1995 to 2021, and no later year is in the archive.
RAD2 measures every channel. RAD1 measures $32$ to $38$ of its channels in nearly every
minute of every year, from $20$ to $1040$\,kHz at spacings that grow from $4$\,kHz at
the bottom to about $100$\,kHz at the top, and the archive fills the other channels by
interpolation between the measured ones, which the files mark through each day's minimum
voltage. STEREO's high-frequency receiver
covers $125$\,kHz to $16$\,MHz in $50$\,kHz channels about every $39$\,s, from CDAWeb's
level-3 files, whose spectra are calibrated in W\,m$^{-2}$\,Hz$^{-1}$ after the
subtraction of a daily background. We take every spectrum of STEREO-A in 2008, 2012 and
2020 and of STEREO-B in 2010 and 2012, years spread over the spacecraft's drift from
Earth and, for 2012, held by both.

\emph{Lines.} In each spectrum a channel carries a line when it stands five robust
deviations above the median of its eight neighbouring channels. A channel's occupancy is
the fraction of a year's spectra in which it carries a line, and a line is persistent
when its occupancy exceeds one per cent on thirty days or more. In RAD1 a measured
channel stands above interpolated neighbours whenever the spectrum curves, so its lines
are tested among the measured channels alone, each against the median of its eight
nearest measured neighbours. A line narrower than a channel is diluted over the channel's
width, so the statistic is a line's power against the background in the channel, and a
line is narrowband here only in the sense of filling one channel of $3$, $20$ or
$50$\,kHz.

\emph{Known sources.} A persistent line is assigned, in order, to the spacecraft's own
lines; to the shortwave broadcast bands within $25$\,kHz; to Earth's auroral kilometric
radiation between $100$ and $700$\,kHz and the solar wind's plasma line below
$100$\,kHz; to the maritime and aeronautical mobile bands within $25$\,kHz and the
amateur bands; to broadcasting up to $200$\,kHz outside the broadcast bands; and to the
terrestrial sources whose occupancy by hour of UT varies by more than its mean, since
Earth's transmitters wax and wane as it turns. The service bands with their $25$\,kHz
margins cover $64$ per cent of RAD2's range and $52$ per cent of STEREO's, and with the
$200$\,kHz of out-of-band broadcasting $76$ and $62$ per cent; the auroral and
plasma-line ranges cover $67$ per cent of RAD1's. A line inside those ranges is
assigned to the service or the natural source and not counted, so the search for lines
outside the known sources covers the remaining quarter of RAD2's band, $38$ per cent of
STEREO's, and the third of RAD1's above $700$\,kHz.

\emph{Sensitivity.} The channel-to-channel scatter of a one-minute spectrum sets what
a line must carry. In RAD2 the robust deviation of a channel from the median of its
eight neighbours is $2.4$ per cent of the background power, in every year sampled
(1999 and 2020), so a $5\sigma$ line adds $12$ per cent of the background to a channel.
Above $1$\,MHz the background is the Galaxy's, whose brightness is close to
$10^{-20}$\,W\,m$^{-2}$\,Hz$^{-1}$\,sr$^{-1}$ from $2$ to $10$\,MHz~\cite{Cane1979},
and a short dipole collects it from $8\pi/3$\,sr, so the background's equivalent flux
density is $8\times10^{-20}$\,W\,m$^{-2}$\,Hz$^{-1}$ and the threshold for a line
filling RAD2's $20$\,kHz is $F=5\times0.024\times8\times10^{-20}\times2\times10^{4}
\approx2\times10^{-16}$\,W\,m$^{-2}$, for a source at the antenna's best direction.
That is the flux of an isotropic source of $6\times10^{7}$\,W at $1$\,au,
$2\times10^{18}$\,W at $1$\,pc and $2\times10^{20}$\,W at $10$\,pc, or of a transmitter
of $5\times10^{7}$\,W behind a kilometre-wide aperture at $1$\,pc, whose beam at
$5$\,MHz is $60$\,rad across and gains nothing. STEREO's calibrated spectra scatter by
$6$ to $7$ per cent of a background of $1.4\times10^{-19}$\,W\,m$^{-2}$\,Hz$^{-1}$
above $2$\,MHz in the raw receiver power and by $23$ per cent of the
background-subtracted level, so its $50$\,kHz channels reach about
$2\times10^{-15}$\,W\,m$^{-2}$. RAD1's measured channels scatter about the median of
their measured neighbours by $3$ per cent of the background in 2020 and by $10$ per
cent above $300$\,kHz and $29$ per cent below it in 1999, where the auroral emission
varies from minute to minute, so a line must exceed its neighbours by $15$ per cent of
the background in the one year and by half to one and a half times it in the other;
with the galactic brightness absorbed to a few
$10^{-21}$\,W\,m$^{-2}$\,Hz$^{-1}$ below $1$\,MHz the threshold in a $3$\,kHz channel
is of order $10^{-17}$ to $10^{-16}$\,W\,m$^{-2}$, where the receiver's own noise,
which dominates the lowest channels, raises it further.

\emph{Day-to-day correlation.} Terrestrial leakage reaches the spacecraft only through a
sunlit ionosphere, so every leaked line rises and falls from day to day with the
broadcast lines of the $41$\,m band ($7225$ to $7375$\,kHz). The Spearman correlation of
a line's daily occupancy with theirs separates leakage from a line made on the
spacecraft or arriving from beyond Earth.

\emph{Spacecraft lines.} A narrowband source beyond the spacecraft reaches Wind and both
STEREO spacecraft, whose receivers are about equally sensitive, about equally often. A
line that one spacecraft holds five times as often as Wind, or as the other STEREO
spacecraft in the same year, is made on that spacecraft.

\emph{Instrument changes.} Two intervals change Wind's spectra. During its spacecraft
anomaly of 2014 October 26 to 2015 January 30 the $1175$\,kHz channel switched on and off,
and in 2014 and 2015 RAD1's list of measured channels shrank to $15$; both are set aside.

% =============================================================================
\section{Results}\label{sec:results}
% =============================================================================

\emph{Wind.} In every year from 1995 to 2021 the persistent lines all have known
sources. In RAD2 there are $119$ persistent lines in 1995, $153$ in 1996, $112$ in 1998
and $153$ in 1999: in 1999, $55$ in the broadcast bands, $36$ in the maritime and $13$ in
the aeronautical bands, $14$ in out-of-band broadcasting, $2$ in the amateur bands, $31$
with the diurnal modulation of a terrestrial source, and $2$ of Wind's own. One line
outside the broadcast bands,
$6975$\,kHz in 1997, follows the $41$\,m lines from day to day (Spearman $0.52$), as
terrestrial leakage must, and so does $7175$\,kHz in 2020 ($0.53$). The steady line at
$8025$\,kHz, present on $365$ and $330$ days of 2020 and 2021 with no daily modulation,
does not ($-0.16$, against $0.74$ for the broadcast line at $9625$\,kHz), so it is made
on Wind; so are the lines at $1175$\,kHz and, in 1995, at $1275$ and $1525$\,kHz. In
RAD1 the test against all neighbouring channels marks $55$ to $96$ lines a year, every
one a measured channel standing above its interpolated neighbours, which is the
interpolation and not a line. Tested among the measured channels alone over 2020, the
$48$ channels RAD1 measured hold $26$ persistent lines, all below $500$\,kHz: eight
below $100$\,kHz, where the solar wind's plasma line wanders, and $18$ from $104$ to
$484$\,kHz in the range of the auroral kilometric radiation, with occupancies of $1.4$
to $7.7$ per cent and the day-to-day variation of that emission; the six measured
channels above $700$\,kHz hold none.

\emph{STEREO.} Table~\ref{tab:stereo} lists the STEREO years. STEREO-A holds $7$, $8$
and $10$ persistent lines in 2008, 2012 and 2020 and STEREO-B $6$ in each of 2010 and
2012, every one in a service band, in the range of auroral kilometric radiation, or made
on the spacecraft. The spacecraft lines include
$1275$, $1475$, $1675$ and $2075$\,kHz, which Wind holds less than a fifth as often,
$725$\,kHz in STEREO-A in 2012, which STEREO-B holds $0.7$ per cent of the time against
STEREO-A's $25$, and $10425$\,kHz, which STEREO-B holds $87$ to $90$ per cent of the time
in 2010 and 2012 and STEREO-A $0.08$, $0.4$ and $4.7$ per cent in 2008, 2012 and 2020,
while Wind's occupancy there stays near $0.7$ per cent: a line of the receiver design
that grew in STEREO-A.

\begin{table}[t]
\centering\small
\caption{The STEREO years. Persistent lines exceed one per cent occupancy on thirty days
or more; left are those outside the known sources. Spacecraft lines are held five times
as often as by Wind or the other STEREO spacecraft in the same year; the others lie in
service bands or in the range of auroral kilometric radiation.}
\label{tab:stereo}
\begin{tabular}{l c r r >{\raggedright\arraybackslash}p{0.45\textwidth}}
\toprule
Spacecraft & Year & Persistent & Left & Spacecraft lines (kHz) \\
\midrule
STEREO-A & 2008 & $7$ & $0$ & $1375$, $1675$ \\
STEREO-A & 2012 & $8$ & $0$ & $725$, $1275$, $1475$, $1675$, $2075$ \\
STEREO-A & 2020 & $10$ & $0$ & $725$, $1275$, $1475$, $1675$, $6925$, $10425$ \\
STEREO-B & 2010 & $6$ & $0$ & $6975$, $10425$, $10475$ \\
STEREO-B & 2012 & $6$ & $0$ & $6975$, $10425$, $10475$ \\
\bottomrule
\end{tabular}
\end{table}

% =============================================================================
\section{Related work}\label{sec:related}
% =============================================================================

\emph{Low-frequency searches.} Searches for narrowband technosignatures at the lowest
frequencies observed from the ground, near $100$\,MHz, have used wide-field arrays such
as the Murchison Widefield Array~\cite{Tingay2016}, whose channels of tens of kilohertz
and sensitivity of tens of janskys over hours make them, like this one, searches for
powerful broadband-channel emitters rather than for hertz-wide lines; we find no search
below the ionospheric cutoff.

\emph{Spacecraft radio receivers.} The receivers of Wind and STEREO were built for solar
radio bursts, interplanetary plasma waves and Earth's radio emissions, and their
calibration against the Galaxy's background~\cite{Dulk2001} makes them sensitive
monitors of the whole low-frequency sky. The Radio Astronomy Explorer first showed how
much terrestrial noise leaks to space~\cite{Herman1973}.

% =============================================================================
\section{Conclusion}\label{sec:conclusion}
% =============================================================================

No persistent narrowband line below the ionospheric cutoff, from $20$\,kHz to $16$\,MHz,
lies outside the known natural, terrestrial and spacecraft sources in $27$ years of Wind
spectra and five years of STEREO spectra. Terrestrial leakage betrays itself by
following the broadcast bands from day to day, and lines made on a spacecraft by being
missing from the others.

The search covers lines present in more than one per cent of the minutes on a month of
days, five robust deviations above the neighbouring channels in a minute's average,
which in RAD2 is $2\times10^{-16}$\,W\,m$^{-2}$ in a $20$\,kHz channel: a transmitter
of $6\times10^{7}$\,W radiating isotropically anywhere within $1$\,au, or of
$2\times10^{18}$\,W at $1$\,pc, so the result bounds emitters in the Solar System and
only the most powerful conceivable ones at the nearest stars, where a beam gains
nothing at these wavelengths. A
transmitter that is brief, drifting or broader than a channel is not constrained, nor is
one inside the bands of terrestrial services, which are set aside whole and cover two
thirds of the band above $1$\,MHz; RAD1 constrains only its $32$ to $38$ measured
channels. Below a few
hundred kilohertz the warm interstellar clouds absorb, so that part of the band bounds
transmitters within the heliosphere and along the shortest paths through warm gas. The STEREO
spacecraft's direction finding would give the arrival direction of any line the
spacecraft share, and the years of STEREO not yet searched would extend the comparison
across the spacecraft's whole drift.

\section*{Acknowledgements}

This work uses Wind/WAVES and STEREO/WAVES data from NASA's Coordinated Data Analysis Web
(CDAWeb) at the Space Physics Data Facility, and the software NumPy, SciPy and cdflib.

\begin{thebibliography}{99}\small\raggedright
\setlength{\itemsep}{1pt}

\bibitem{Tingay2016}
S.~J. Tingay, C.~Tremblay, A.~Walsh and R.~Urquhart, An opportunistic search for
extraterrestrial intelligence (SETI) with the Murchison Widefield Array,
\emph{Astrophys. J. Lett.} \textbf{827}, L22 (2016).

\bibitem{RybickiLightman1979}
G.~B. Rybicki and A.~P. Lightman, \emph{Radiative Processes in Astrophysics} (Wiley, New
York, 1979).

\bibitem{RedfieldLinsky2008}
S.~Redfield and J.~L. Linsky, The structure of the local interstellar medium. IV.
Dynamics, morphology, physical properties, and implications of cloud-cloud
interactions, \emph{Astrophys. J.} \textbf{673}, 283 (2008).

\bibitem{Bougeret1995}
J.-L. Bougeret, M.~L. Kaiser, P.~J. Kellogg, et al., WAVES: the radio and plasma wave
investigation on the Wind spacecraft, \emph{Space Sci. Rev.} \textbf{71}, 231 (1995).

\bibitem{Bougeret2008}
J.-L. Bougeret, K.~Goetz, M.~L. Kaiser, et al., S/WAVES: the radio and plasma wave
investigation on the STEREO mission, \emph{Space Sci. Rev.} \textbf{136}, 487 (2008).

\bibitem{Gurnett1974}
D.~A. Gurnett, The Earth as a radio source: terrestrial kilometric radiation, \emph{J.
Geophys. Res.} \textbf{79}, 4227 (1974).

\bibitem{Herman1973}
J.~R. Herman, J.~A. Caruso and R.~G. Stone, Radio Astronomy Explorer (RAE)-I
observations of terrestrial radio noise, \emph{Planet. Space Sci.} \textbf{21}, 443
(1973).

\bibitem{Dulk2001}
G.~A. Dulk, W.~C. Erickson, R.~Manning and J.-L. Bougeret, Calibration of low-frequency
radio telescopes using the galactic background radiation, \emph{Astron. Astrophys.}
\textbf{365}, 294 (2001).

\bibitem{Cane1979}
H.~V. Cane, Spectra of the non-thermal radio radiation from the galactic polar regions,
\emph{Mon. Not. R. Astron. Soc.} \textbf{189}, 465 (1979).

\end{thebibliography}

\end{document}
